\chapter{Clasificación y parentesco}\label{ch:g6:classification-kinship}

¿Por qué funciona tan bien la \omterm{def:g4:classifying-animals:classification}{clasificación} del
\cref{ch:g4:classifying-animals}? Clasifica \omterm{def:g1:animals-around-us:animal}{animales} por las \omterm{def:g1:animals-around-us:covering}{plumas} y
el \omterm{def:g1:animals-around-us:covering}{pelo} y otros mil rasgos se ponen en fila obedientemente --- los que
alimentan con \omterm{def:g2:animals-grow-up:born}{leche} tienen \omterm{def:g4:classifying-animals:vertebrate}{columna vertebral}, los que llevan \omterm{def:g1:animals-around-us:covering}{plumas}
ponen \omterm{def:g2:animals-grow-up:born}{huevos} con cáscara, y nadie tiene que comprobarlo. La
\cref{rem:g4:classifying-animals:tree} dejó la pregunta en el aire.
Este capítulo la responde con una sola palabra --- \omterm{prop:g6:classification-kinship:kinship}{parentesco} --- y
reconstruye las cajas como un \omterm{def:g6:classification-kinship:tree}{árbol} de familia.

\section{Atributos compartidos y encajados}

\begin{definition}[Atributo]\label{def:g6:classification-kinship:attribute}
Un \emph{atributo}\index{atributo} es un rasgo que un organismo
\emph{posee}, enunciado de manera que se pueda comprobar: tiene un
\omterm{def:g3:bones-and-muscles:skeleton}{esqueleto} con \omterm{def:g4:classifying-animals:vertebrate}{columna vertebral}; tiene cuatro \omterm{def:g1:my-body:parts}{extremidades}; tiene \omterm{def:g1:animals-around-us:covering}{pelo}
y alimenta con \omterm{def:g2:animals-grow-up:born}{leche}; tiene \omterm{def:g1:animals-around-us:covering}{plumas}. La \omterm{def:g4:classifying-animals:classification}{clasificación} se construye sobre
atributos --- los buenos \omterm{def:g3:sorting-living-things:criterion}{criterios} de la
\cref{def:g3:sorting-living-things:criterion}, formalizados.
\end{definition}

\begin{proposition}[Los atributos encajan unos dentro de otros]\label{prop:g6:classification-kinship:nest}
Comprobados en muchas \omterm{def:g5:biodiversity:species}{especies}, los \omterm{def:g6:classification-kinship:attribute}{atributos} caen en conjuntos
\emph{encajados} --- todos los \omterm{def:g1:animals-around-us:animal}{animales} con \omterm{def:g1:animals-around-us:covering}{pelo} tienen además cuatro
\omterm{def:g1:my-body:parts}{extremidades}; todos los \omterm{def:g1:animals-around-us:animal}{animales} de cuatro \omterm{def:g1:my-body:parts}{extremidades} tienen además
\omterm{def:g4:classifying-animals:vertebrate}{columna vertebral}; nunca al revés (\omterm{def:g1:animals-around-us:covering}{pelo} sin \omterm{def:g4:classifying-animals:vertebrate}{columna vertebral}) y nunca
cruzándose (\omterm{def:g1:animals-around-us:covering}{pelo} y \omterm{def:g1:animals-around-us:covering}{plumas} a la vez). Ese encaje es la razón de que los
grupos dentro de grupos
(\cref{def:g4:classifying-animals:classification}) le sirvan de algo al
mundo vivo.
\end{proposition}
\admitted

\begin{omfigure}
\begin{tikzpicture}[scale=0.9]
  % nested attribute boxes
  \draw[thick, omDef, fill=omDef!4] (0,0) rectangle (11.4,4.4);
  \node[font=\small\bfseries, omDef] at (2.2,4.1)
    {tiene columna vertebral};
  \draw[thick, omMeth, fill=omMeth!5] (0.4,0.35) rectangle (8.6,3.5);
  \node[font=\small\bfseries, omMeth] at (2.35,3.2)
    {tiene cuatro extremidades};
  \draw[thick, omProp, fill=omProp!6] (0.8,0.7) rectangle (5.6,2.6);
  \node[font=\small\bfseries, omProp] at (2.6,2.3)
    {tiene pelo y leche};
  \node[font=\small] at (3.2,1.4) {gato, murciélago, ballena};
  \node[font=\small, align=center] at (7.0,1.6)
    {rana, lagartija,\\ mirlo};
  \node[font=\small] at (10.0,1.6) {trucha};
\end{tikzpicture}

{\small Los \omterm{def:g6:classification-kinship:attribute}{atributos} encajan unos dentro de otros: los que llevan \omterm{def:g1:animals-around-us:covering}{pelo}
están dentro de los de cuatro \omterm{def:g1:my-body:parts}{extremidades}, y estos dentro de los que
tienen \omterm{def:g4:classifying-animals:vertebrate}{columna vertebral}. La trucha tiene la columna, pero se queda
fuera de la caja de las cuatro \omterm{def:g1:my-body:parts}{extremidades}.}
\end{omfigure}

\begin{example}[Colocar tres animales]\label{ex:g6:classification-kinship:placing}
Comprueba las cajas para la ballena, la trucha y el mirlo. Ballena:
\omterm{def:g4:classifying-animals:vertebrate}{columna vertebral} sí, cuatro \omterm{def:g1:my-body:parts}{extremidades} sí (aletas delanteras y
vestigios escondidos detrás), \omterm{def:g1:animals-around-us:covering}{pelo} y \omterm{def:g2:animals-grow-up:born}{leche} sí --- la caja más interior,
con el gato. Trucha: \omterm{def:g4:classifying-animals:vertebrate}{columna vertebral} sí, y se acabó --- solo la caja
exterior. Mirlo: \omterm{def:g4:classifying-animals:vertebrate}{columna vertebral}, cuatro \omterm{def:g1:my-body:parts}{extremidades} (dos patas, dos
alas), \omterm{def:g1:animals-around-us:covering}{plumas} --- la caja de las cuatro \omterm{def:g1:my-body:parts}{extremidades}, en el
compartimento propio de las \omterm{def:g1:animals-around-us:covering}{plumas}. Ningún \omterm{def:g1:animals-around-us:animal}{animal} exige nunca dos cajas
cruzadas a la vez.
\end{example}

\section{El parentesco: la razón que hay detrás del encaje}

\begin{proposition}[Los atributos compartidos significan antepasados compartidos]\label{prop:g6:classification-kinship:kinship}
La explicación del encaje es la herencia (la
\cref{rem:g4:animal-reproduction:mix}: las crías reciben de sus
progenitores). Las \omterm{def:g5:biodiversity:species}{especies} que comparten un \omterm{def:g6:classification-kinship:attribute}{atributo} lo comparten
porque lo \emph{heredaron de un antepasado común} --- una \omterm{def:g5:biodiversity:species}{especie}
abuela lejana. Cuantos más \omterm{def:g6:classification-kinship:attribute}{atributos} comparten dos \omterm{def:g5:biodiversity:species}{especies}, más
cercano es su \emph{parentesco}\index{parentesco}: el gato y la
ballena, los dos peludos y con \omterm{def:g2:animals-grow-up:born}{leche}, son primos más cercanos entre sí
de lo que cualquiera de los dos lo es de la trucha. Un grupo de una
\omterm{def:g4:classifying-animals:classification}{clasificación} es, en el fondo, una \emph{familia}: un antepasado y todos
sus descendientes.
\end{proposition}
\admitted

\begin{example}[Leer por fin bien a la ballena]\label{ex:g6:classification-kinship:whale}
Por eso el rompecabezas de la ballena
(\cref{exo:g4:classifying-animals:9}) tenía que resolverse como se
resolvió: el \omterm{def:g1:animals-around-us:covering}{pelo} y la \omterm{def:g2:animals-grow-up:born}{leche} marcan la línea familiar de la ballena ---
sus antepasados eran \omterm{prop:g3:sorting-living-things:groups}{mamíferos} terrestres que alimentaban con \omterm{def:g2:animals-grow-up:born}{leche} ---
mientras que su forma de \omterm{prop:g3:sorting-living-things:groups}{pez} es cosa del agua
(\cref{exo:g4:adapted-to-their-world:11}), estampada sobre una familia
sin \omterm{prop:g6:classification-kinship:kinship}{parentesco} por una misma manera de vivir. Los \omterm{def:g6:classification-kinship:attribute}{atributos} de la
ascendencia clasifican familias; los parecidos de \omterm{def:g2:where-animals-live:habitat}{hábitat} se los pueden
prestar unos desconocidos. La \omterm{def:g4:classifying-animals:classification}{clasificación} se fía de los primeros y
desconfía de los segundos.
\end{example}

\begin{definition}[El árbol del parentesco]\label{def:g6:classification-kinship:tree}
El \omterm{prop:g6:classification-kinship:kinship}{parentesco} se dibuja como un \emph{árbol}\index{árbol del
parentesco}: los puntos de ramificación representan antepasados comunes
y cada \omterm{def:g6:classification-kinship:attribute}{atributo} aparece una sola vez, en la \omterm{def:g1:plants-around-us:tree}{rama} donde surgió --- todas
las \omterm{def:g5:biodiversity:species}{especies} que están más allá de esa bifurcación lo heredan. Las
cajas encajadas y el árbol llevan la misma información: una caja es una
\omterm{def:g1:plants-around-us:tree}{rama} con todas sus ramitas.
\end{definition}

\begin{omfigure}
\begin{tikzpicture}[scale=0.95]
  % tree
  \draw[very thick] (0,0) -- (2.2,0);
  % backbone mark
  \fill[omDef] (1.6,0) circle (0.09);
  \node[font=\footnotesize, omDef, align=center] at (1.6,-0.55)
    {aparece la\\ columna vertebral};
  \draw[very thick] (2.2,0) -- (3.6,1.6) -- (10.4,1.6);
  \draw[very thick] (2.2,0) -- (4.4,-1.4) -- (10.4,-1.4);
  \node[font=\small, right] at (10.5,-1.4) {trucha};
  % four limbs mark
  \fill[omMeth] (3.1,1.05) circle (0.09);
  \node[font=\footnotesize, omMeth, align=center] at (2.4,1.45)
    {aparecen las\\ cuatro extremidades};
  \draw[very thick] (5.0,1.6) -- (6.2,2.8) -- (10.4,2.8);
  % fur mark on upper branch
  \fill[omProp] (5.75,2.35) circle (0.09);
  \node[font=\footnotesize, omProp, align=center] at (5.0,2.9)
    {aparecen el pelo\\ y la leche};
  \draw[very thick] (6.9,2.8) -- (7.7,3.7) -- (10.4,3.7);
  \node[font=\small, right] at (10.5,3.7) {gato};
  \node[font=\small, right] at (10.5,2.8) {ballena};
  % feathers on middle branch
  \fill[omThm] (7.4,1.6) circle (0.09);
  \node[font=\footnotesize, omThm, align=center] at (7.4,1.05)
    {aparecen las\\ plumas};
  \node[font=\small, right] at (10.5,1.6) {mirlo};
\end{tikzpicture}

{\small Un \omterm{def:g6:classification-kinship:tree}{árbol} de \omterm{prop:g6:classification-kinship:kinship}{parentesco} de bolsillo. Cada punto de color es un
\omterm{def:g6:classification-kinship:attribute}{atributo} que aparece una sola vez; todas las \omterm{def:g1:plants-around-us:tree}{ramas} que hay más allá lo
heredan. El gato y la ballena comparten la bifurcación más reciente ---
los primos más cercanos que se dibujan aquí.}
\end{omfigure}

\begin{method}[Construir un árbol de parentesco pequeño]\label{met:g6:classification-kinship:build}
Dado un puñado de \omterm{def:g5:biodiversity:species}{especies}:
\begin{enumerate}
  \item haz una tabla con sus \omterm{def:g6:classification-kinship:attribute}{atributos} --- las \omterm{def:g5:biodiversity:species}{especies} en filas, los
        \omterm{def:g6:classification-kinship:attribute}{atributos} comprobables en columnas y una marca donde se posean;
  \item agrupa por marcas compartidas, empezando por el \omterm{def:g6:classification-kinship:attribute}{atributo} más
        amplio: el encaje aparece en la tabla;
  \item dibuja el \omterm{def:g6:classification-kinship:tree}{árbol}: una bifurcación por \omterm{def:g6:classification-kinship:attribute}{atributo}, la más amplia
        más cerca del tronco; coloca cada \omterm{def:g5:biodiversity:species}{especie} pasando exactamente
        las bifurcaciones cuyos \omterm{def:g6:classification-kinship:attribute}{atributos} posee;
  \item lee el \omterm{prop:g6:classification-kinship:kinship}{parentesco}: la cercanía de dos \omterm{def:g5:biodiversity:species}{especies} es la última
        bifurcación que comparten.
\end{enumerate}
\end{method}

\begin{remark}[Lo que el árbol da a entender sin decirlo]\label{rem:g6:classification-kinship:implies}
Tómate en serio el \omterm{def:g6:classification-kinship:tree}{árbol} y estará diciendo algo enorme: sigue hacia
atrás dos \omterm{def:g1:plants-around-us:tree}{ramas} cualesquiera y \emph{se juntan}. El gato y la trucha
comparten una \omterm{def:g5:biodiversity:species}{especie} abuela; así que, más atrás todavía, tienen que
compartirla también el gato y el caracol, la margarita y un niño. Todos
los \omterm{def:g6:exploring-our-environment:components}{seres vivos} serían entonces una sola familia, y su unidad
(\cref{prop:g5:first-look-at-cells:principle}) un aire de familia --- y
las \omterm{def:g5:biodiversity:species}{especies} tendrían que poder \emph{cambiar} de una generación a
otra, o la ramificación no habría podido ocurrir nunca. Si eso es
cierto, y cómo podría funcionar, es la gran pregunta del curso del
\cref{ch:g9:evidence-of-evolution} --- este \omterm{def:g6:classification-kinship:tree}{árbol} es el dibujo que la
obliga a plantearse.
\end{remark}

\section{Ejercicios}

\begin{exercise}[$\star$]\label{exo:g6:classification-kinship:1}
¿Qué es un \omterm{def:g6:classification-kinship:attribute}{atributo}? Da tres, enunciados de manera comprobable.
\end{exercise}

\begin{exercise}[$\star$]\label{exo:g6:classification-kinship:2}
¿Qué significa que los \omterm{def:g6:classification-kinship:attribute}{atributos} \emph{encajen} unos dentro de otros?
Da el ejemplo de las tres cajas del capítulo.
\end{exercise}

\begin{exercise}[$\star$]\label{exo:g6:classification-kinship:3}
¿Por qué comparten \omterm{def:g6:classification-kinship:attribute}{atributos} las \omterm{def:g5:biodiversity:species}{especies}, según la
\cref{prop:g6:classification-kinship:kinship}?
\end{exercise}

\begin{exercise}[$\star$]\label{exo:g6:classification-kinship:4}
En el \omterm{def:g6:classification-kinship:tree}{árbol} de bolsillo, ¿qué punto de \omterm{def:g6:classification-kinship:attribute}{atributo} tiene detrás la trucha?
¿Qué puntos tiene detrás la ballena?
\end{exercise}

\begin{exercise}[$\star$]\label{exo:g6:classification-kinship:5}
¿Quiénes son primos más cercanos: la ballena y el gato o la ballena y
la trucha? Lee la respuesta en las bifurcaciones del \omterm{def:g6:classification-kinship:tree}{árbol}.
\end{exercise}

\begin{exercise}[$\star$]\label{exo:g6:classification-kinship:6}
Enuncia la relación entre una caja encajada y una \omterm{def:g1:plants-around-us:tree}{rama} del \omterm{def:g6:classification-kinship:tree}{árbol}.
\end{exercise}

\begin{exercise}[$\star\star$]\label{exo:g6:classification-kinship:7}
¿Por qué se fía la \omterm{def:g4:classifying-animals:classification}{clasificación} del \omterm{def:g1:animals-around-us:covering}{pelo} y la \omterm{def:g2:animals-grow-up:born}{leche} pero desconfía de
``tiene forma de \omterm{prop:g3:sorting-living-things:groups}{pez}''? Usa el \cref{ex:g6:classification-kinship:whale}.
\end{exercise}

\begin{exercise}[$\star\star$]\label{exo:g6:classification-kinship:8}
Construye la tabla de \omterm{def:g6:classification-kinship:attribute}{atributos} del
\cref{met:g6:classification-kinship:build} para: trucha, rana, gato y
mirlo --- \omterm{def:g6:classification-kinship:attribute}{atributos}: \omterm{def:g4:classifying-animals:vertebrate}{columna vertebral}, cuatro \omterm{def:g1:my-body:parts}{extremidades}, \omterm{def:g1:animals-around-us:covering}{pelo} y
\omterm{def:g2:animals-grow-up:born}{leche}, \omterm{def:g1:animals-around-us:covering}{plumas}.
\end{exercise}

\begin{exercise}[$\star\star$]\label{exo:g6:classification-kinship:9}
Añade la lagartija al \omterm{def:g6:classification-kinship:tree}{árbol} de bolsillo: ¿qué bifurcaciones pasa y por
dónde sale su \omterm{def:g1:plants-around-us:tree}{rama}?
\end{exercise}

\begin{exercise}[$\star\star$]\label{exo:g6:classification-kinship:10}
El murciélago vuela como un mirlo pero amamanta como un gato. ¿Qué
parecido marca el \omterm{prop:g6:classification-kinship:kinship}{parentesco} y cuál está prestado por la manera de
vivir? Decídelo con la lógica del \omterm{def:g6:classification-kinship:tree}{árbol}.
\end{exercise}

\begin{exercise}[$\star\star$]\label{exo:g6:classification-kinship:11}
``Un grupo es un antepasado y todos sus descendientes''. Pon a prueba
la frase con los \omterm{prop:g3:sorting-living-things:groups}{mamíferos} y con el ornitorrinco del
\cref{exo:g4:classifying-animals:11}: ¿rompe la familia esa rareza que
pone \omterm{def:g2:animals-grow-up:born}{huevos} o pertenece a ella?
\end{exercise}

\begin{exercise}[$\star\star\star$]\label{exo:g6:classification-kinship:12}
Lleva la \cref{rem:g6:classification-kinship:implies} hasta su límite:
¿qué tiene que ser cierto de las \omterm{def:g5:biodiversity:species}{especies} a lo largo del tiempo si
todos los pares de \omterm{def:g1:plants-around-us:tree}{ramas} se juntan de verdad? Enuncia la implicación en
una frase --- y nombra qué clase de pruebas exigirías para ella.
\end{exercise}

\section{Problema: El cajón del museo}

\begin{problem}[{Problema de fin de semana --- seis ejemplares, una
tabla de atributos y un árbol}]\label{pb:g6:classification-kinship:1}
Un museo le presta a la clase un cajón con seis ejemplares: una trucha,
una rana, una lagartija, una paloma, un ratón y un \omterm{def:g3:bones-and-muscles:skeleton}{cráneo} de delfín con
moldes de sus aletas. La tarea: clasificarlos por \omterm{def:g6:classification-kinship:attribute}{atributos} y dibujar
su \omterm{def:g6:classification-kinship:tree}{árbol} de \omterm{prop:g6:classification-kinship:kinship}{parentesco}.

\medskip\noindent\textbf{Parte I --- La tabla.}
\begin{enumerate}
  \item Construye la tabla de los seis, con los \omterm{def:g6:classification-kinship:attribute}{atributos}: \omterm{def:g4:classifying-animals:vertebrate}{columna
        vertebral}; cuatro \omterm{def:g1:my-body:parts}{extremidades}; \omterm{def:g1:the-five-senses:sense}{piel} seca y escamosa; \omterm{def:g1:animals-around-us:covering}{plumas};
        \omterm{def:g1:animals-around-us:covering}{pelo} y \omterm{def:g2:animals-grow-up:born}{leche}. (Para el delfín, casi sin \omterm{def:g1:animals-around-us:covering}{pelo}, la etiqueta del
        museo da fe: crías alimentadas con \omterm{def:g2:animals-grow-up:born}{leche} --- el
        \cref{ex:g3:sorting-living-things:surprises} explica por qué
        ese testimonio lo zanja.)
  \item ¿Qué \omterm{def:g6:classification-kinship:attribute}{atributo} comparten los seis? ¿Qué significa, en términos
        de \omterm{prop:g6:classification-kinship:kinship}{parentesco}, que se comparta tan ampliamente?
  \item ¿Qué ejemplares marcan ``cuatro \omterm{def:g1:my-body:parts}{extremidades}''? Explica las
        marcas de la paloma y del delfín --- ¿qué son las alas y las
        aletas, según la lógica del \omterm{def:g6:classification-kinship:tree}{árbol}?
  \item ¿Por qué no aparece ``vive en el agua'' en ninguna parte de la
        tabla, si la trucha y el delfín viven los dos en ella? Cita la
        regla del capítulo.
\end{enumerate}

\medskip\noindent\textbf{Parte II --- El \omterm{def:g6:classification-kinship:tree}{árbol}.}
\begin{enumerate}[resume]
  \item Dibuja el \omterm{def:g6:classification-kinship:tree}{árbol}: ordena las bifurcaciones --- \omterm{def:g4:classifying-animals:vertebrate}{columna
        vertebral}, cuatro \omterm{def:g1:my-body:parts}{extremidades} y después las cubiertas
        (escamas, \omterm{def:g1:animals-around-us:covering}{plumas}, \omterm{def:g1:animals-around-us:covering}{pelo}) en sus \omterm{def:g1:plants-around-us:tree}{ramas} propias --- y coloca a los
        seis.
  \item Lee en tu \omterm{def:g6:classification-kinship:tree}{árbol}: el primo más cercano del ratón dentro del
        cajón y el de la trucha.
  \item La rana está más allá de ``cuatro \omterm{def:g1:my-body:parts}{extremidades}'' pero antes de
        todas las bifurcaciones de cubierta. ¿Qué la marca su \omterm{def:g1:the-five-senses:sense}{piel}
        desnuda y húmeda
        (\cref{rem:g4:classifying-animals:confusion}) y por dónde sale
        su \omterm{def:g1:plants-around-us:tree}{rama}?
  \item A un séptimo ejemplar se le cae la etiqueta: \omterm{def:g1:the-five-senses:sense}{piel} seca y
        escamosa, cuatro \omterm{def:g1:my-body:parts}{extremidades}, \omterm{def:g4:classifying-animals:vertebrate}{columna vertebral}. Colócalo en
        el \omterm{def:g6:classification-kinship:tree}{árbol} y nombra su clase --- sin ver la etiqueta.
\end{enumerate}

\medskip\noindent\textbf{Parte III --- El significado.}
\begin{enumerate}[resume]
  \item El delfín y la trucha comparten el cuerpo afinado; el delfín y
        el ratón comparten la ascendencia con \omterm{def:g1:animals-around-us:covering}{pelo} y \omterm{def:g2:animals-grow-up:born}{leche}. ¿Qué pareja
        está más cerca en el \omterm{def:g6:classification-kinship:tree}{árbol} y por qué tiene el \omterm{def:g6:classification-kinship:tree}{árbol} que ignorar
        el cuerpo afinado?
  \item Enuncia qué afirma históricamente la bifurcación de las
        ``cuatro \omterm{def:g1:my-body:parts}{extremidades}'': ¿de qué única cosa, hace mucho
        tiempo, descienden la rana, la lagartija, la paloma, el ratón y
        el delfín?
  \item El \omterm{def:g6:classification-kinship:tree}{árbol} de la clase hace que todas las \omterm{def:g1:plants-around-us:tree}{ramas} se junten, tarde
        o temprano, en el tronco. Di en una frase qué afirma sin
        decirlo ese dibujo sobre los seis ejemplares --- y sobre los
        alumnos que lo dibujan.
  \item Un escéptico dice: ``tu \omterm{def:g6:classification-kinship:tree}{árbol} no es más que un sistema de
        archivo''. Da la respuesta en dos partes: ¿de qué es un
        \emph{hecho} el encaje de los \omterm{def:g6:classification-kinship:attribute}{atributos} y qué \emph{explica} el
        \omterm{def:g6:classification-kinship:tree}{árbol} que un simple archivo dejaría sin explicar?
\end{enumerate}
\end{problem}
